\section{Consequences and comparison with existing theory}\label{sec:scope}
The acquired law, its geometric resolution and the propagation of implementation error are three different objects. They are coupled here by a common experiment and task, not by an ambient parameter dimension. Theorem~\ref{thm:main} makes this coupling quantitative. Theorem~\ref{thm:kernel} supplies its causal estimate from one-step marked-kernel data; the normalized matrix in \eqref{eq:resolvent-gamma} records, rather than suppresses, the remaining comparison with acquired mass. Theorem~\ref{thm:robust} then gives a single fully specified experiment class in which acquisition, memory and compulsory information losses occur together.

The three realizations test distinct parts of this mechanism. The filter derives a posterior density from the actual report Jacobian. The absorbing sensor derives a mixed-dimensional acquired measure and a disjoint-indicator transport identity. The recurrent experiment derives a density bound under repeated folding and refresh, and has active expansions even though its forced second moment is bounded. Each follows the general theorem; none is a premise of it.

\subsection{Theorem-level comparison}
\paragraph{Predictive states and sufficient statistics.}
Shalizi and Crutchfield's minimality and uniqueness theorems for causal states \cite[Theorems 2--3]{shalizi} identify minimal prescient representations of a stochastic process. Predictive state representations \cite{psr} use future action--observation tests. Our quotient proposition does not claim priority for either construction. It explicitly includes controlled continuation legality, paid resources and conditional-version domains. Its extra measurable work is the density-instrument descent at every declared continuous report. The categorical treatment of kernels and sufficiency in \cite{fritz} supplies another general language; an abstract sufficient statistic alone does not provide the acquired-mass or finite-machine estimates used here.

\paragraph{Bisimulation and belief-state control.}
Behavioral metrics for Markov decision processes \cite{ferns} compare states through rewards and transition laws and control value differences. Our metric is fixed by executable prediction tasks and valid update geometry; it is not asserted to be a universal bisimulation metric. Saldi, Y\"uksel and Linder \cite[Assumption 1 and Theorem 3]{saldi} obtain near-optimal finite belief-model policies for discounted partially observed control under their continuity, compactness and moment hypotheses. That theorem optimizes a controller, whereas our matching lower is for a fixed exploration and the law it actually acquires. Our finite-horizon two-sided resolution estimate does not imply their unrestricted control conclusion, nor does asymptotic policy near-optimality supply our actual-mass lower bound.

\paragraph{Quantization and nonlinear filtering.}
Conditional-mean projection and optimal probability quantization are classical \cite{graf}. Nonlinear filter approximation by quantization is developed in \cite{pages,sellami,feuer}. The additional obligation here is a matching lower for the very same acquired task, together with the forced-moment condition needed by a persistent finite machine. Zhu's local estimates for Ahlfors--David measures \cite[Theorem 1.1]{zhu} concern a substantially broader singular-measure setting than a finite union of box charts. Our finite-budget anisotropic bounds expose constants through width collapse and intersecting finite mixtures; they neither classify arbitrary singular measures nor replace established quantization-dimension theory.

\paragraph{Positive moment operators.}
Average contraction of iterated random functions \cite{diaconis} and moment stability of Markov jump systems \cite{costa} are established subjects. We claim no novelty for a Neumann resolvent or the criterion $\rho(K_2)<1$. The role of Theorem~\ref{thm:kernel} is more specific: reverse-edge operators propagate the conditional first and mixed second moments of the \emph{rounding forcing}, and the resulting matrix is normalized by the actual acquisition weights. Exact suspension cancels in the forced profiles, not in the physical budget. This distinction is why a recurrent expanding kernel can still yield a matching finite-memory risk curve.

\paragraph{Comparison of experiments.}
Blackwell comparison \cite{blackwell}, Le Cam's deficiency theory \cite{lecam}, and Torgersen's systematic account \cite{torgersen} provide the classical common-decision-loss principle. The triangle inequality and total-variation contraction in Theorem~\ref{thm:morphism} are instances of that principle. The added structure is the finite-resource causal implementation, with explicit lifted controller classes, stopping indices, clocks and numerical interfaces. A forward simulator yields an upper risk transfer, not a lower bound relative to a different oracle. The mandatory-state and robust converses are proved directly in Theorems~\ref{thm:tag} and \ref{thm:robust}, rather than attributed to a forward deficiency bound.

\subsection{Necessary distinctions illustrated by exact examples}
First, the valid envelope $[0,1]$ equipped with a point mass and the same envelope equipped with a uniform law have respectively zero one-point distortion and strictly positive finite-label distortion. Hence a reachability envelope is not an acquired law. More generally a component of probability $w$ contributes with weight $w$, even if its geometric diameter is large.

Second, an identity update has no need for repeated quantization. If the initial state has already been replaced by a representative, the same representative can be retained indefinitely. The bound $Nr$ obtained by blindly injecting a covering radius at every hold is not an inherent resource requirement. The acquisition-weighted certificate corrects precisely this failure of accounting.

Third, a statistical garbling can increase memory needed to match its \emph{own} oracle. Observing a fair hidden bit exactly produces only posterior probabilities zero and one. Adding independent Gaussian noise to this observation produces a continuously varying posterior probability. Two labels reproduce the first oracle exactly, but no finite number reproduces the latter continuously distributed oracle exactly under square score. This is compatible with data processing because the baselines differ.

Fourth, let $P_N=\delta_0$ and $Q_N=(1-e^{-Nc})\delta_0+e^{-Nc}\delta_1$, $c>0$. Their total variation distance tends to zero, but the exponential rates at the point one are infinite and $c$, respectively. The common-risk defect of this paper is therefore not a certificate of large-deviation equivalence.

Finally, a hidden mode or an internal simulator register is not free because its name is omitted from a state vector. Independent pairs of retained states have the product cardinality in Theorem~\ref{thm:morphism}. Conversely that product is a constructive upper bound, not a theorem that every simulator requires it. Theorem~\ref{thm:tag} establishes necessity only through separately executable, zero-error continuation challenges. An exact label machine and a finite-bit numerical implementation are also different objects: the latter retains the nonzero calibration, workspace, and readout terms of \eqref{eq:joint}. Exact-arithmetic zero distortion is not a promise of a finite-precision physical output with zero error.


\subsection{Uniformity and the remaining mathematical questions}
The conclusions apply to the experiment classes actually specified. For the structural route, a marked factor, invariant environment law, dominating preparation, bounded Lyapunov weight and acquired-weight normalization must be verified. In the one-chart recurrent class the normalization is automatic, and active expansion is permitted. In many-chart classes it can fail when a terminal component becomes rare while earlier innovations still transport into it. A sharper finite-horizon localization is then a mathematical problem, not a free consequence of the resolvent bound.

The total-resource converse is nonconstant in the mandatory continuation memory. It does not establish a shortest-program, minimum-workspace or optimal-runtime converse below the constructive feasible region, and a small allowed tag-error probability would require a quantitative relaxation of the continuation-separation argument. The robust precision floor comes from compulsory indistinguishable input cells; optional numerical inaccuracy has no such converse. These are different proof responsibilities even though all coordinates are retained in the upper resource map.

The framework also leaves concrete extensions to be proved: learning unknown kernels with optimally charged pilots; unrestricted adaptive exploration with a common lower-bound experiment; infinitely accumulating strata beyond the finite-chart mass bounds; and adaptive or state-dependent suspensions not covered by Proposition~\ref{prop:suspension}. Bounded stopping comparisons are proved, but an infinite-horizon transcript comparison, a large-deviation equivalence or closure for unbounded generators would require stronger and different estimates. None is inferred from finite-time total variation or from a realization-specific identity. The central program remains the same experiment-to-geometry-to-resource transfer, rather than a relabeling of any single realization.
